\let\Bbbk\undefined
\usepackage{amsthm,proof,amssymb}
\usepackage{enumitem}
\usepackage{oktheorem,okgeneralmath,okcategory}
\usepackage{bbold}
\usepackage{epiolmec}
\usepackage{oksemantics}
\usepackage{mathpartir}
\usepackage{todonotes}
\usepackage{natbib}
\usepackage{fp}
\usepackage{xparse}
\usepackage{xspace}

%next two for no-depth underlining
\usepackage{contour}
%\usepackage[normalem]{ulem}
\newcommand{\OhK}[1]{\todo[inline,author=OK]{#1}}
\newcommand{\ohk}[1]{\todo[size=\tiny]{OK: #1}{}}
\newcommand{\MaP}[1]{\todo[inline,author=MP,color=LightGreen]{#1}}
\newcommand{\map}[1]{\todo[size=\tiny,color=LightGreen]{MP: #1}{}}


%%%%%%%%%%% Cross references %%%%%%%%%%
\newboolean{WithSupplementary}
\setboolean{WithSupplementary}{true}
\newcommand\seclabel[1]{\label{sec:#1}}
\newcommand\secref[1]{\S\ref{sec:#1}}
\newcommand\applabel[1]{\label{app:#1}}
\newcommand\appref[2][]{Appendix~%
  \ifthenelse{\boolean{WithSupplementary}}%
             {\ref{app:#2}}%
             {#1}%
             }
\newcommand\applemmaref[2][]{%
  \ifthenelse{\boolean{WithSupplementary}}%
             {\lemmaref{#2}}%
             {Lemma~#1}%
}
\newcommand\apptheoremref[2][]{%
  \ifthenelse{\boolean{WithSupplementary}}%
             {\theoremref{#2}}%
             {Theorem~#1}%
}
\WithSuffix\newcommand\appref*[1]{\ref{app:#1}}
\newcommand\subseclabel[1]{\label{subsec:#1}}
\newcommand\subsecref[1]{\S\ref{subsec:#1}}
\newcommand\subsubseclabel[1]{\label{subsubsec:#1}}
\newcommand\remarklabel[1]{\label{remark:#1}}
\newcommand\remarkref[1]{Remark~\ref{remark:#1}}
\newcommand\subsubsecref[1]{\S\ref{subsubsec:#1}}
\newcommand\figlabel[1]{\label{fig:#1}}
\newcommand\figref[1]{Fig.~\ref{fig:#1}}

\newcommand\figfile[2]{%
    \begin{figure}%
      \begin{center}%
        \input{sections/#1}%
        \caption{#2}%
        \figlabel{#1}%
      \end{center}%
    \end{figure}%
}

{
  \theoremstyle{definition}
  \newtheorem{remark}[definition]{Remark}
}

%%%%%%%%%%%%%%%%%%%%%%%%%%%%%%%%%%%%%%


%%%%%%%%%%% Add possessive citations, e.g. "Leroy's (1992)", taken
\newcommand{\citepos}[1]{\citeauthor{#1}'s~\cite{#1}}
\WithSuffix\newcommand\citepos*[2]{\citeauthor{#1}'s #2~\cite{#1}}
\WithSuffix\newcommand\citepos-[2]{\citeauthor{#1}#2~\cite{#1}}
%%%%%%%%%%%%%%%% End of possessive citation %%%%%%%%%%%%%%%%%%%%%
\makeatletter
\newcommand\ibidt[2][]{%
  \citeauthor{#2}~\ibid[#1]{#2}
}
\newcommand\ibid[2][]{%
  \citetext{%
    \hyper@natlinkstart{#2}%
    \textit{ibid.}%
    \hyper@natlinkend%
    \ifx\ibid@empty#1\ibid@empty%
    \else%
    , #1%
    \fi%
  }%
}

\newcommand\citloct[1]{%
  \citeauthor{#1}~\citloc{#1}%
}
\newcommand\citloc[1]{%
  \citetext{%
    \hyper@natlinkstart{#1}%
    \textit{cit.~loc.}%
    \hyper@natlinkend%
  }%
}
\newcommand\citlocpos[1]{%
  \citeauthor{#1}'s~\citloc{#1}%
}
\WithSuffix\newcommand\citlocpos*[2]{%
  \citeauthor{#1}'s #2~\citloc{#1}%
}
\makeatother


\newcommand\diagram[1]{\begin{center}
    \diagram*{#1}
\end{center}}
\WithSuffix\newcommand\diagram*[2][c]{\begin{tabular}[#1]{@{}c@{}}\includegraphics{generated-diagrams/diagram-#2.mps}\end{tabular}}
\newcommand\diagramNumbered[2][]{%
  \begin{equation}
  #1
  \hfill
  \diagram*{#2}
  \hfill
  \end{equation}%
}

%%%%%%%%%%%%%%%%%%%%%%%%%%%%%%%%%%%%%%
\usepackage{tcolorbox}
\newtcbox{\myfigbox}{%
  on line, arc=3pt, outer arc=3pt,
  colback=black!8, colframe=black!8,
  boxsep=0pt,
  left=2pt,right=2pt,top=1pt,bottom=1pt,%
  boxrule=0pt}
%%%%%%%%%%%%%%%%%%%%%%%%%%%%%%%%%%%%%
% Import a single character from a font to conserve font tables
\makeatletter
\NewDocumentCommand\NewFontImport{O{\hbox}O{\tf@size}mmmmmm}{
\newcommand#3{\begingroup\mathchoice%
  {#1{\fontsize{#2}{#2}\usefont{#4}{#5}{#6}{#7}\char#8}}%
  {\hbox{\fontsize{\tf@size}{\tf@size}\usefont{#4}{#5}{#6}{#7}\char#8}}%
  {\hbox{\fontsize{\sf@size}{\sf@size}\usefont{#4}{#5}{#6}{#7}\char#8}}%
  {\hbox{\fontsize{\ssf@size}{\ssf@size}\usefont{#4}{#5}{#6}{#7}\char#8}}%
  \endgroup%
}}
\makeatother

%%%%%%%%%%%%%%%%%%%%%%%%%%%%%%%%%%%%%%
\makeatletter
\newcommand\newmvar[3][]{%
  \newcommand#2[1][1]{#1{\ifcase ##1\or #3\or\else@ctrerr\fi}}%
}

%% introduce a variation of a defined meta-variable
\newcommand\mvarvar[3]{%
  \WithSuffix\newcommand#1#2[1][0]{%
    \FPeval{\tempindex}{clip(#3+##1)}%
    #1[\tempindex]%
  }
}

\makeatother


%% half-open intervals, since they mess up indentation
\newcommand\oc[2]{(#1, #2]}
\newcommand\co[2]{[#1, #2)}

\newcommand\symdiff{\mathbin{\triangle}}

%%%%%%%%%%%%%%%%%%%%%%%%%%%%%%%%%%%
%%%%%%%%%%% monads %%%%%%%%%%%%%
\newcommand\newop[2]{\newcommand{#1}{\mathop{{\mathrm{#2}}}\nolimits}}
\newcommand{\bindsymbol}{\mathrel{\small\rrangle}\joinrel\mathrel{\small=}}
\newcommand{\klseq}{\ggg}
\newcommand{\bind}{\mathrel{\bindsymbol}}
\newop{\return}{return}
\newop{\strengthName}{\mathbf{st}}
\newcommand{\strength}[1][]{\supimplicit[{#1}]{\strengthName}}
\newcommand{\xstrength}[1][]{\supimplicit[{#1};\times]{\strengthName}}
\newop{\shadow}{shadow}

\newcommand\suspend[1]{\boldsymbol\{#1\boldsymbol\}}
\newcommand\spair[2]{\langle #1, #2\rangle}
\newcommand\sunit{\langle\rangle}
\makeatletter
\let\llangle\@undefined
\let\rrangle\@undefined
%% \DeclareMathDelimiter{\llangle}{\mathopen}%
%%                      {okMnSymbolLargeSymbols}{'164}{okMnSymbolLargeSymbols}{'164}
%% \DeclareMathDelimiter{\rrangle}{\mathclose}%
%%                      {okMnSymbolLargeSymbols}{'171}{okMnSymbolLargeSymbols}{'171}
\makeatother
\def\pexplen{0mu}
\def\pexpmidlen{12mu}
\newcommand\pexp{
  \mathrel{%
    \overrightharpoon{%
      {\mspace{-\pexplen}\smash{\mathord{\relbar\mspace{-\pexpmidlen}\relbar}}\vphantom{.}}\mspace{-\pexplen}}%
  }%
}
\usepackage{mathtools}
\usepackage{pst-venn}
\usepackage{tikz-cd}
\usetikzlibrary{decorations.pathmorphing}

% code from Gonzalo
\newcommand\xrsquigarrow[1]{%
    \mathrel{%
        \begin{tikzpicture}[%
            baseline={(current bounding box.south)}
            ]
        \node[%
            ,inner sep=.44ex
            ,align=center
            ] (tmp) {$\scriptstyle #1$};
        \path[%
            ,draw,<-
            ,decorate,decoration={%
            ,snake,aspect=0
            ,amplitude=1pt
            ,segment length=1.5mm,pre length=3.5pt
                }
            ]
        (tmp.south east) -- (tmp.south west);
        \end{tikzpicture}
        }
    }
\newcommand\xlsquigarrow[1]{%
    \mathrel{%
        \begin{tikzpicture}[%
            ,baseline={(current bounding box.south)}
            ]
        \node[%
            ,inner sep=.44ex
            ,align=center
            ] (tmp) {$\scriptstyle #1$};
        \path[%
            ,draw,<-
            ,decorate,decoration={%
                ,zigzag
                ,amplitude=0.7pt
                ,segment length=1.2mm,pre length=3.5pt
                }
            ]
        (tmp.south west) -- (tmp.south east);
        \end{tikzpicture}
        }
    }

%%%%%%%%%%%%%%%%%%%%%%%%%%%%%%%%%%%%%%%%%%
\usepackage[dvipsnames,svgnames,table]{xcolor}
\usepackage{listings}
\colorlet{code}{DarkBlue}
\colorlet{types}{Magenta}
\colorlet{implicit}{Gray}
%%%%%%%%%%%%%%%%% Syntax %%%%%%%%%%%%%%%%%
\newcommand\gdefinedby{\mathrel{::=}}
\newcommand\gor{\lvert}


\newcommand\syncat[1]{\mspace{-25mu}\synname{#1}}
\newcommand\synname[1]{\qquad\text{#1}}

\newenvironment{syntax}[1][]{%
\(
  \rowcolors{100}{white}{white}%
  \begin{array}[t]{#1l@{}l@{\,}*2{l@{}}@{\,}l}
}{
\end{array}
\)%
}

\newenvironment{sugar}{%
  \(
  \begin{array}[t]{@{\bullet\ }l@{\mspace{20mu}}ll}
    \multicolumn{1}{@{\hphantom{\bullet}}l}{\text{Sugar}} & \text{Elaboration}\\
    \hline
}{
  \end{array}
  \)%
}

\usepackage{scalerel,stackengine}
\stackMath
\newcommand\reallywidehat[1]{%
\savestack{\tmpbox}{\stretchto{%
  \scaleto{%
    \scalerel*[\widthof{\ensuremath{#1}}]{\kern-.6pt\bigwedge\kern-.6pt}%
    {\rule[-\textheight/2]{1ex}{\textheight}}%WIDTH-LIMITED BIG WEDGE
  }{\textheight}%
}{0.5ex}}%
\stackon[1pt]{#1}{\tmpbox}%
}

\newcommand\sem[1]{\scottbrackets{#1}}

\newenvironment{catdef}{%
\begingroup%
\newcommand\objects[1]{\item[Objects $##1$:]}%
\newcommand\morphisms[1]{\item[{Morphisms $##1$:}]}%
\begin{description}}{\end{description}\endgroup}

\NewDocumentCommand\eval{oO{}}{\mathop{\mathrm{eval}}\nolimits_{\IfNoValueF{#1}{#1;}#2}}
\NewDocumentCommand\curry{oO{}}{\mathop{\mathrm{curry}}\nolimits_{\IfNoValueF{#1}{#1;}#2}}
\NewDocumentCommand\uncurry{oO{}}{\mathop{\mathrm{uncurry}}\nolimits_{\IfNoValueF{#1}{#1;}#2}}
\RenewDocumentCommand\Dom{oO{}}{\mathop{\mathrm{Dom}}\nolimits_{\IfNoValueF{#1}{#1;}#2}}
\newcommand\swap{\mathtt{F}}
\newcommand\stet{\mathtt{I}}
\newcommand\pto\rightharpoonup
\newcommand\xpto\xrightharpoonup

\newcommand\genepto[2]{
\mathrel{\raisebox{#1}{$#2\lhook$}}\joinrel\rightleftharpoons
}

\newcommand\epto{
  \mathchoice{\genepto{1pt}{\displaystyle}}
             {\genepto{1pt}{\textstyle}}
             {\genepto{0.8pt}{\scriptstyle}}
             {\genepto{0.8pt}{\scriptscriptstyle\mkern -1mu}}
}

\newcommand\xeptogen[2]{
\mathrel{\raisebox{#1}{$#2\lhook$}}\joinrel\xrightleftharpoons
}

\newcommand\xepto[2][]{
  \mathchoice{\xeptogen{.9pt}{\displaystyle}[#1]{#2}}
             {\xeptogen{1pt}{\textstyle}[#1]{#2}}
             {\xeptogen{0.8pt}{\scriptstyle}[#1]{#2}}
             {\xeptogen{0.8pt}{\scriptscriptstyle\mkern -1mu}[#1]{#2}}
}
\newcommand\dif{\mathrm{d}}
\newcommand\Dif{\mathop\triangle\!}
\DeclareFontFamily{U}{bskarr}{}
\DeclareFontShape{U}{bskarr}{m}{n}{<->bskarr10}{ }
\NewFontImport\rightarrowTriangle{U}{bskarr}{m}{n}{"FB}
\NewFontImport\rightarrowtriangle{U}{bskarr}{m}{n}{"FE}
\NewDocumentCommand\kto{O{}}{\mathrel{\rightarrowtriangle}_{%
  \ifthenelse{\equal{#1}{sig}}%
          {\sigma}{%
  \ifthenelse{\equal{#1}{prb}}{1}{%
  \ifthenelse{\equal{#1}{fin}}{\textrm{fin}}{%
  \ifthenelse{\equal{#1}{sub}}{\leq1}
             {#1}}}}%
}}

\DeclareFontEncoding{LS1}{}{}
\DeclareFontFamily{LS1}{stixarrows1}{ }
\DeclareFontSubstitution{LS1}{stix}{m}{n}
\DeclareFontShape{LS1}{stixarrows1}{m}{n} {<-> stix-mathsf}{}
%\NewFontImport\lhook{LS1}{stixarrows1}{m}{n}{"2A}
\newcommand{\xbothooks}{%
  \mathchoice{\mathrel{\lhook}\mspace{-4mu}\joinrel\mathrel{\raisebox{-2.6pt}{$\displaystyle\lhook$}}}
             {\mathrel{\lhook}\mspace{-3mu}\joinrel\mathrel{\raisebox{-2.6pt}{$\lhook$}}}
             {\mathrel{\raisebox{.6pt}{$\scriptstyle\lhook$}}%
                 \mspace{-3mu}\joinrel%
              \mathrel{\raisebox{-1.05pt}{$\scriptstyle\lhook$}}}
             {\mathrel{\raisebox{1.5pt}{$\scriptscriptstyle\lhook$}}%
                 \mspace{-2mu}\joinrel%
              \mathrel{\raisebox{-.19pt}{$\scriptscriptstyle\lhook$}}}}
\newcommand{\bothooks}{%
  \mathchoice{\mathrel{\lhook}\mspace{-3.5mu}\joinrel\mathrel{\raisebox{-2.6pt}{$\displaystyle\lhook$}}}
             {\mathrel{\lhook}\mspace{-3mu}\joinrel\mathrel{\raisebox{-2.6pt}{$\lhook$}}}
             {\mathrel{\lhook}%
                 \mspace{-3mu}\joinrel%
              \mathrel{\raisebox{-1.9pt}{$\scriptstyle\lhook$}}}
             {\mathrel{\raisebox{.2pt}{$\scriptscriptstyle\lhook$}}%
                 \mspace{-2mu}\joinrel%
              \mathrel{\raisebox{-1.25pt}{$\scriptscriptstyle\lhook$}}}}
\DeclareFontFamily{LS1}{stixintegrals}{ }
\DeclareFontShape{LS1}{stixintegrals}{m}{n} {<-> stix-mathcal}{}
%\NewFontImport[\raisebox{-2pt}][15pt]\sqintslop{LS1}{stixintegrals}{m}{n}{"A8}
\NewFontImport\intcupupop{LS1}{stixintegrals}{m}{n}{"C7}
\newcommand\intcupup{\DOTSI\mathop{\intcupupop}\limits}
\DeclareFontFamily{LS1}{stixletters}{ }
\DeclareFontShape{LS1}{stixletters}{m}{it}{<-> stix-mathit}{}
\NewFontImport{\scurel}{LS1}{stixletters}{m}{it}{"FE}
\DeclareFontFamily{LS1}{stixfrak}{ }
\DeclareFontShape{LS1}{stixfrak}{m}{n}{<-> stix-mathfrak}{}
\NewFontImport{\boxastSymb}{LS1}{stixfrak}{m}{n}{"B5}
\newcommand\mapmul{\mathbin{\boxastSymb}}
\makeatletter
%\newcommand\sqint{\DOTSI\mathop{\sqintslop}\nolimits}
\makeatother

\newcommand\cotuple\coseq
\newcommand\deneq{\mathrel{%
  \raisebox{.5ex}{%
  \scalebox{0.4}{\EOx}%
  }
  }}
\newcommand\vto[3]{%
  \begin{smallmatrix}%
        \hphantom{#2}#1     \\%
        #2\mathord\downarrow\\%
        \hphantom{#2}#3       %
  \end{smallmatrix}%
}
\WithSuffix\newcommand\vto+[3]{%
  \begin{smallmatrix}%
        #1     \\%
        \mathllap{#2}\mathord\downarrow\\%
        #3       %
  \end{smallmatrix}%
}
\WithSuffix\newcommand\vto*[3]{%
  \begin{smallmatrix}%
        #1\hphantom{#2}     \\%
        \mathord\downarrow#2\\%
        #3\hphantom{#2}       %
  \end{smallmatrix}%
}
\WithSuffix\newcommand\vto-[3]{%
  \begin{smallmatrix}%
        #1     \\%
        \mathord\downarrow\mathrlap{#2}\\%
        #3       %
  \end{smallmatrix}%
}
\newcommand\vtotwo[5]{%
  \begin{smallmatrix}%
        \hphantom{#2}#1     \\%
        #2\mathord\downarrow\\%
        \hphantom{#2}#3     \\%
        #4\mathord\downarrow\\%
        \hphantom{#4}#5       %
  \end{smallmatrix}%
}
\WithSuffix\newcommand\vtotwo+[5]{%
  \begin{smallmatrix}%
        #1                             \\%
        \mathllap{#2}\mathord\downarrow\\%
        #3                             \\%
        \mathllap{#4}\mathord\downarrow\\%
        #5
  \end{smallmatrix}%
}
\WithSuffix\newcommand\vtotwo*[5]{%
  \begin{smallmatrix}%
        #1\hphantom{#2}     \\%
        \mathord\downarrow#2\\%
        #3\hphantom{#2}     \\%
        \mathord\downarrow#4\\%
        #5\hphantom{#4}       %
  \end{smallmatrix}%
}
\newcommand\vfrom[3]{%
  \begin{smallmatrix}%
        \hphantom{#2}#1     \\%
        #2\mathord\uparrow\\%
        \hphantom{#2}#3       %
  \end{smallmatrix}%
}

\newcommand\para{\mathbin{\mathrel\lvert\mathrel\lvert}}

%% \DeclareFontFamily{U}{mathb}{}
%% \DeclareFontShape{U}{mathb}{m}{n}{
%%       <5> <6> <7> <8> <9> <10> gen * mathb
%%       <10.95> mathb10 <12> <14.4> <17.28> <20.74> <24.88> mathb12
%%       }{}
%% \DeclareFontSubstitution{U}{bbold}{m}{n}
%%\NewFontImport\llcorner{mathb}{"7A}
%%\NewFontImport\lrcorner{mathb}{"7B}


%\newcommand\carrier[2][]{\llcorner\underset{#1}{#2}\lrcorner}
\newcommand\carrier[1]{\underline{#1}}
\newcommand\free[2][]{\ulcorner\overset{#1}{#2}\urcorner}
\newcommand\fun[1]{#1 \mapsto}
\newenvironment{casefun}{\uplambda\begin{cases}}{\end{cases}}
\newcommand\mathi{\mathchoice%
  {\text{\normalsize\rm\i}}%
  {\text{\normalsize\rm\i}}%
  {\text{\scriptsize\rm\i}}%
  {\text{\tiny\rm\i}}%
}

\newenvironment{codeblock}%
   {\begin{array}[t]{@{}l@{}}}
   {\end{array}}

\DeclareFontFamily{U}{bbold}{}
\DeclareFontShape{U}{bbold}{m}{n}{<->bbold10}{}
\DeclareFontSubstitution{U}{bbold}{m}{n}
\NewFontImport\SinSym{U}{bbold}{m}{n}{'017}

%%%%%%%%%%%%%  Metavariables %%%%%%%%%%%%%%%%%%%%%%%%%%%%%%

% algebras
\newmvar\A{A \or B \or C \or D}
\newcommand\B{\A[2]}
\let\C\undefined
\newcommand\C{\A[3]}
\newcommand\D{\A[4]}
% Sets
\newmvar\X{X \or Y \or Z}
\newcommand\Y{\X[2]}
\newcommand\Z{\X[3]}
\newmvar\mX{X \or Y \or Z}
\newcommand\mY{\mX[2]}
\newcommand\mZ{\mX[3]}
\newmvar\mx{x \or y \or z}
\newcommand\my{\mx[2]}
\newcommand\mz{\mx[3]}

\newmvar\E{E}
% Groups
\newmvar\G{G \or H \or K}
\newmvar\M{M \or K \or L}
\let\H\undefined
\newcommand\H{\G[2]}
\newcommand\K{\G[3]}
% Predicates
\let\P\undefined
\newmvar\P{P \or Q}
\newcommand\Q{\P[2]}
% Representations
\newmvar\rep{\rho\or\sigma\or\tau}
% Actions are algebras

\newcommand\implicit[2][]{\overset{{\color{implicit}#1}}{#2}}
% inlined implicit
\newcommand\inplicit[1]{{\color{implicit}#1}}
\newcommand\subimplicit[2][]{#2_{{\color{implicit}#1}}}
\newcommand\supimplicit[2][]{#2^{{\color{implicit}#1}}}
\newcommand\Group{\Category{Grp}}
\newcommand\presentation[2]{\langle #1 \mid #2 \rangle}
\newcommand\unit[1][]{\implicit[{#1}]{e}}
\newcommand\mul[1][]{\mathbin{\implicit[#1]{\ast}}}
\newcommand\ginv[2][]{\subimplicit[#1]{\inv{#2}}}
\WithSuffix\newcommand\ginv*[2][]{\subimplicit[#1]{\inv*{#2}}}
\newcommand\Mul[1][]{\mathop{\implicit[#1]{\bigast}}}
\newcommand\Action[1]{#1\text{-}\Category{Act}}
\WithSuffix\newcommand\Action*{\Category{Act}}
\newcommand\FaithfulAction{\Category{FaithfulAct}}
\newcommand\ConcGroup{\Category{ConcreteGrp}}
\newcommand\ConcAction{\Category{ConcreteAct}}
\newcommand\srAction[1]{#1\text{-}\Category{Act}_{\mathrm{sr}}}
\newcommand\act[1][]{\mathbin{{\circledast}^{\inplicit{#1}}}}
\WithSuffix\newcommand\act'[1][]{\mathop{\implicit[#1]{\odot}}\nolimits}
\newcommand\lact[2][\act]{#1#2}
\newcommand\ract[2][\act]{#2#1}
\newcommand\iact[2][]{\mathbin{{\circledast}^{\inplicit{#1}}_{#2}}}
\WithSuffix\newcommand\iact*[2][]{\mathbin{{\circledast}^{\inplicit{#1}_{\inplicit{#2}}}}}
\WithSuffix\newcommand\iact+[2][]{\act[#1]}
\newcommand\Orb[1][\,]{\mathsf{Orb}#1}
\newcommand\Triv[1][]{\mathsf{Triv}_{#1}}
\newcommand\Trans[1][\,]{\mathsf{Trans}#1}
\newcommand\adjunit[1][]{\implicit[#1]{\eta}}
\newcommand\adjcounit[1][]{\implicit[#1]{\epsilon}}
\newcommand\wlg{\mathop{\mathsf{wlog}}\nolimits}
\newcommand\equito{\mathrel{\raisebox{.25ex}{$\scriptscriptstyle\circledast$}}\joinrel\to}
\newcommand\homoto{\mathrel{\raisebox{.15ex}{$\scriptstyle\ast$}}\joinrel\to}
\newcommand\equifrom{\from\joinrel\mathrel{\raisebox{.25ex}{$\scriptscriptstyle\circledast$}}}
\newcommand\xequifrom[1]{\overset{#1}\equifrom}
\newcommand\xequito[1]{\overset{#1}\equito}

\newcommand\types\vdash
\newcommand\flatten{\mathop{\mathsf{flatten}}\nolimits}

\newcommand\Ker{\mathop{\mathrm{Ker}}\nolimits}
\newcommand\WKer{\mathop{{\mathrm{\tilde{K}er}}}\nolimits}
\newcommand\Canon[2][]{\mathop{\mathsf{Can}}\nolimits_{#1}{#2}}
\newcommand\canoniser{c}
\newcommand\SemiRegAction[1]{\Action{#1}_{\mathrm{sreg}}}

\newcommand\normsub{\triangleleft}
\newcommand\groth\oint
\newcommand\quotact{\mathop{\mathsf{quotAct}}\nolimits}
\newcommand\stabilise[2][\G]{#1_{#2}}
\newcommand\stabeq[1]{\overset{#1}\sim}
\newcommand\Fin{\mathop{\mathsf{\TConst{Fin}}}\nolimits}

\newcommand\Sym[1]{\mathbf{S}_{#1}}
\newcommand\Cyc[1]{\mathbf{C}_{#1}}
\newcommand\Dih[1]{\mathbf{D}_{#1}}
\newcommand\Trn[1]{\mathbf{T}_{#1}}
\newcommand\Dag[1]{\mathbf{DAG}_{#1}}
\newcommand\Graph[1]{\mathbf{Graph}_{#1}}
\newcommand\TopDag[1]{\mathbf{DAG}^\mathrm{top}_{#1}}
\newmvar\perm{\pi\or\sigma\or\tau}
\newcommand\TopSort{\mathop{\mathsf{TopSort}}\nolimits}
\newcommand\topSort{\mathop{\mathsf{topSort}}\nolimits}
\newcommand\topSortDepthFirst{\mathop{\mathsf{topSortDepthFirst}}\nolimits}
\newcommand\sort{\mathop{\mathsf{sort}}\nolimits}
\newcommand\insrt{\mathop{\mathsf{insert}}\nolimits}
\newcommand\Natinf{\naturals_{\infty}}
\newcommand\dist{\mathop{\mathsf{dist}}\nolimits}
\newcommand\Dist[1]{\mathbf{Dist}_{#1}}
\newcommand\Prop{\mathbf{\TConst{Prop}}}
\newmvar\prop{\phi\or\psi\or\xi}
\newcommand\mprop[1][1]{\BoundV{\prop[#1]}}
\newcommand\length\cardinality
\newcommand\Bool{\mathbf{Bool}}
\WithSuffix\newcommand\++{\mathbin{+\mspace{-8mu}{+}}}
\WithSuffix\newcommand\h+{\mathbin{{\color{red}+}\mspace{-8mu}{\color{blue}+}}}
\newcommand\preleq{\sqsubseteq}
\newcommand\pregeq{\sqsupseteq}
\newcommand\prelt{\sqsubset}
\newcommand\pregt{\sqsupset}
\newcommand\LL[1]{{#1}_{\infty}}
\newcommand\Total[1][\keys]{\mathop{\mathbf{Total}}_{#1}}
\newcommand\Bound[1]{\mathop{\mathbf{Bound}}_#1}
\newcommand\Interval[1]{\mathop{\mathsf{Interval}}#1}
\newcommand\IsInterval[2][]{\mathop{\mathsf{IsInterval}}_{#1}#2}
\newcommand\Ordering{\mathop{\mathsf{Ordering}}}
\makeatletter
\long\def\ok@tree[#1][#2]{\mathop{\mathsf{T}}_{#1}^{#2}}
\newcommand\SearchTree{\optparams{\ok@tree}{[{}][{}]}}
\makeatother
\WithSuffix\newcommand\SearchTree*[1]{#1\text{-}\SearchTree}
%\newcommand\Tree[1][l \preleq u]{\mathop{\mathbf{Tree}}_{#1}}

\newmvar\btree{t\or s\or r}
\newcommand\mtree[1][1]{\BoundV{\btree[#1]}}
%\newcommand\Node[4][]{\mathsf{Node}\,\inplicit{#1}#2\, #3\, #4}
\newcommand\NodeName{\mathsf{Node}}
\WithSuffix\newcommand\leq?[1][\preleq]{\overset{?}{#1}}
\WithSuffix\newcommand\in?{\mathrel{\overset?\in}}
\WithSuffix\newcommand\=?{\mathrel{\overset?=}}
\WithSuffix\newcommand\in!{\mathrel{\overset!\in}}
\WithSuffix\newcommand\notin!{\mathrel{\overset!{\not\in}}}
\newcommand\Aye{\mathsf{Yes}}
\newcommand\Nay{\mathsf{No}}
\newcommand\May{\mathsf{M}}
\newcommand\At{\mathbin\mathbf{@}}
\newmvar\Ctx{\Gamma\or\Delta\or\Xi\or\Psi\or\Phi}
\newcommand\Dtx{\Ctx[2]}
\newmvar\ctx{\gamma\or\delta\or\xi\or\psi\or\varphi}
\newcommand\dtx{\ctx[2]}
\newcommand\xtx{\ctx[3]}
\newcommand\With{\mathrel{\mathbf{with}}}
\newcommand\inLeft{\mathop{\mathsf{inLeft}}\nolimits}
\newcommand\inRoot{\mathop{\mathsf{inRoot}}\nolimits}
\newcommand\inRigt{\mathop{\mathsf{inRight}}\nolimits}
\newcommand\prf{\text{prf}}
\newcommand\punct[1]{#1_{*}}
\newcommand\Dec{\mathsf{Dec}}
\newcommand\dec{\mathsf{dec}}

\newcommand\preals{\mathbb{R}_+}
\newcommand\pnaturals{\mathbb{N}_+}

\newcommand\rulename\textsc
\newcommand\NeutRule{\rulename{neutral}}
\newcommand\CompRule{\rulename{compose}}
\newcommand\InvRule{\rulename{inverse}}
\newcommand\Gen[1]{\left\langle #1\right\rangle}
\newcommand\unitval\star
\newcommand\Board[1]{\TConst{Board}_{#1}}
\newcommand\symmap[1]{\subimplicit[\mathsf{symm}]{#1}}
\newcommand\elemap[1]{\subimplicit[\mathsf{elem}]{#1}}
\newcommand\fmap[1][]{\mathop{\subimplicit[{#1}]{F}}}
\newcommand\strings{\mathbb{String}}
\newcommand\chars{\mathbb{Char}}
\newcommand\toUpper{\mathsf{toUpper}}
\newcommand\MkString[1]{{\color{DarkRed}\text{``#1''}}}
\newcommand\MkChar[1]{{\color{DarkRed}\text{`#1'}}}
\newcommand\subst[2]{#1\bracks{#2}}
\newcommand\For[2]{#1{\mathrel{\color{DarkRed}\mapsto}}#2}
\newcommand\rotate[1][\,]{\mathsf{rotate}#1}
\newcommand\norm[1]{\left\lVert#1\right\rVert}
\newcommand\cnorm[2][]{\norm{#2}_{#1}}
\newcommand\Vertex[1][]{\mathsf{V}_{#1}}
\newcommand\Edge[1][]{\mathsf{E}_{#1}}
\newmvar\Grph{G\or H\or K}
\newcommand\walkto\leadsto
\newenvironment{permutation}{%
  \left(%
  \begin{array}{@{}>{\scriptstyle}l@{\scriptstyle{}\mapstio{}}>{\scriptstyle}l@{}}%
  }{
  \end{array}%
  \right)%
  }
\newcommand\Row[1]{\TConst{Row}_{#1}}
\newcommand\Vct[2]{\TConst{Vect}_{#1}#2}
\newcommand\Mat[3]{\TConst{Mat}_{#1\times#2}#3}
\DeclareSymbolFont{okMnSymbolA}{U}{MnSymbolA}{m}{n}
\SetSymbolFont{okMnSymbolA}{bold}{U}{MnSymbolA}{b}{n}
\DeclareFontShape{U}{MnSymbolA}{b}{n}{
    <-6>  MnSymbolA-Bold5
   <6-7>  MnSymbolA-Bold6
   <7-8>  MnSymbolA-Bold7
   <8-9>  MnSymbolA-Bold8
   <9-10> MnSymbolA-Bold9
  <10-12> MnSymbolA-Bold10
  <12->   MnSymbolA-Bold12}{}
\DeclareMathSymbol{\rcurvearrownw}{\mathrel}{okMnSymbolA}{197}
\DeclareMathSymbol{\rcurvearrowleft}{\mathrel}{okMnSymbolA}{194}
\DeclareMathSymbol{\rcurvearrowright}{\mathrel}{okMnSymbolA}{192}
\DeclareMathSymbol{\lcurvearrowsw}{\mathrel}{okMnSymbolA}{190}
\DeclareMathSymbol{\lcurvearrowse}{\mathrel}{okMnSymbolA}{191}
\DeclareMathSymbol{\rcirclearrowright}{\mathrel}{okMnSymbolA}{248}
\DeclareMathSymbol{\downrsquigarrow}{\mathrel}{okMnSymbolA}{171}
\newcommand\Rotation{\mathord{\DConst{\boldsymbol\rcurvearrownw}}}
\newcommand\twiceRotation{\mathord{\DConst{\boldsymbol\rcurvearrowright}}}
\newcommand\unRotation{\mathord{\DConst{\boldsymbol\lcurvearrowse}}}
\newcommand\nonRotation{\mathord{\DConst{\boldsymbol\rcirclearrowright}}}
\newcommand\vplusstyle[3]{%
\mathbin{
  \begin{array}[c]{@{}>{#3}c@{}}
    \color{red}+\\[#1]\color{blue}+\\[#2]
  \end{array}
}}
\WithSuffix\newcommand\v+{%
\mathchoice%
        {\vplusstyle{-2.25ex}{.02ex}{\displaystyle}}%
        {\vplusstyle{-2.25ex}{.02ex}{\textstyle}}%
        {\vplusstyle{-2.25ex}{.2ex}{\scriptstyle}}%
        {\vplusstyle{-2.5ex}{.5ex}{\scriptscriptstyle}}%
}
\newcommand\Tile{\TConst{Tile}}
\newcommand\bimap{\mathsf{map}_2}
\newcommand\transpose{\mathsf{transpose}}
\newcommand\reverse{\mathsf{reverse}}
\newcommand\Direction{\TConst{Dir}}
\newcommand\Rt{\DConst{\mathord\Rightarrow}}
\newcommand\Lt{\DConst{\mathord\Leftarrow}}
\newcommand\Up{\DConst{\mathord\Uparrow}}
\newcommand\Dn{\DConst{\mathord\Downarrow}}
\newcommand\mkLt[1][\,]{\mathsf{mkLft}#1}
\newcommand\rotRow[1][\,]{\mathsf{rotRow}#1}
\newcommand\To\Rightarrow
\newcommand\mergeLeft[1][]{\merge_{\Lt}^{#1}}
\newcommand\rw{\mathrm{row}}
\newcommand\merge{\mathsf{merge}}
\newcommand\minimal{\mathsf{indexMin}}
\newcommand\struct[1][]{\subimplicit[{#1}]{F}}
\usepackage[svgnames]{xcolor}
\usepackage{listings}
\colorlet{types}{DarkCyan}
\colorlet{code}{DarkRed}
\colorlet{bound}{MediumPurple}
\newcommand\TConst[1]{\mathsf{\textcolor{types}{#1}}}
\newcommand\DConst[1]{\mathsf{\textcolor{code}{#1}}}
% \newcommand\BoundV[1]{\mathsf{\textcolor{bound}{#1}}}
\newcommand\BoundV[1]{#1}
\newcommand\Nil{\DConst{Nil}}
\newcommand\Leaf{\DConst{Leaf}}
\WithSuffix\newcommand\::{\mathbin{\DConst{::}}}
\newcommand\Node{\mathop{\DConst{Node}}\nolimits}
\newcommand\List{\mathop{\TConst{List}}}
\newcommand\Maybe{\mathop{\TConst{Maybe}}}
\newcommand\Nothing{{\DConst{Nothing}}}
\newcommand\Just{\mathop{\DConst{Just}}}
\newcommand\Tree{\mathop{\TConst{Tree}}}
\newcommand\Data{\mathrel{\mathbf{data}}}
\newcommand\Or{\mathrel{\boldsymbol\vert}}
\newcommand\roll[1][]{\subimplicit[{#1}]{\mathop{\mathbf{roll}}\nolimits}}
\newcommand\unroll[1][]{\subimplicit[{#1}]{\mathop{\mathbf{unroll}}\nolimits}}
\newcommand\vecmap[1][]{\subimplicit[{#1}]{\mathop{\mathbf{map}}\nolimits}}
\newcommand\fold[1][]{\subimplicit[{#1}]{\mathop{\mathbf{fold}}\nolimits}}
\newcommand\pfold[1][]{\subimplicit[{#1}]{\mathop{\mathbf{pfold}}\nolimits}}
\newcommand\Param[1][]{\mathbf{param}_{#1}}
\newcommand\Sort[1][]{\mathbf{sort}_{#1}}
\newcommand\fix{\boldsymbol\mu}
\newcommand\TreeLabel{{\mathbf{Step}}}
\newcommand\TreePath{\TConst{Path}}
\newcommand\Lft{\DConst{L}}
\newcommand\Rgt{\DConst{R}}
\newenvironment{tree diagram}%
  {\scriptsize\begin{tikzcd}[ampersand replacement=\&,column sep=tiny,row sep=tiny, scale=0.3]}%
  {\end{tikzcd}}
\newcommand\treeLink[4]{\arrow["", from=#1-#2, to=#3-#4, no head, no tail]}
\newcommand\NodeV[3]{
 \begin{tree diagram}
      \& #2 \& \\
    \text{$#1$}\&    \& \text{$#3$}
   \treeLink1221
   \treeLink1223
   \end{tree diagram}
 }
\newcommand\NodeL[2]{
 \begin{tree diagram}
      \& #2 \\
    \text{$#1$}\&
   \arrow["", from=1-2, to=2-1, no head, no tail]
 \end{tree diagram}
 }
\newcommand\NodeR[2]{
 \begin{tree diagram}
   #1 \& \\
      \& \text{$#2$}
   \arrow["", from=1-1, to=2-2, no head, no tail]
 \end{tree diagram}
 }

\NewDocumentCommand\gstrength{O{}o}{
  \IfNoValueTF{#2}
    {\supimplicit[{#1}]{\strengthName}}
    {\supimplicit[{#1};{#2}]{\strengthName}}%
 }

\newcommand\WLOG{w.l.o.g.\xspace}
\newcommand\wLOG{W.l.o.g.\xspace}
\newcommand\lookup{\mathsf{lookup}}
\newcommand\lookupElim{\mathsf{lookup}^{\mathrm{Elim}}}
\newcommand\lookupCore{\mathsf{lookup}^{\mathrm{Core}}}
\newcommand\lookLeft{\mathsf{lookLeft}}
\newcommand\PositionIn{\mathsf{\TConst{Pos}}}
\newcommand\NodeInName{\mathsf{\TConst{InPath}}}
\newcommand\NodeIn{\mathop{\NodeInName}}
\newcommand\Type[1]{\mathop{\mathsf{\TConst{#1}}}}
\newmvar\msubst{\theta\or \xi\or\zeta}
\newcommand\OrbitOf[2][\G]{#1\act\bracks{#2}}
\newcommand\weaken[1][]{\mathsf{wkn}_{#1}}
\newmvar\x{\BoundV x\or\BoundV y \or\BoundV z}
\newcommand\y{\x[2]}
\newcommand\z{\x[3]}
\newcommand\DPair[2]{\parent{#1 : #2}\mathbin{\times}}
\newcommand\dpair[1]{#1\mathbin{\DConst,}}
\newcommand\DFun[2]{\parent{#1 : #2}\mathbin{\to}}
\newcommand\EFun[2]{\parent{#1 : #2}\mathbin{\equito}}
\newcommand\Fun[1]{#1\mathbin{\to}}
\newcommand\Tru{\DConst{True}}
\newcommand\Fls{\DConst{False}}
\newmvar\mn{\BoundV m \or \BoundV n\or \BoundV k}
\newmvar\mm{\mn[2]}
\newmvar\mk{\mn[3]}
\newmvar\mt{\BoundV t}
\newmvar\mell{\BoundV \ell}
\newmvar\mr{\BoundV r}
\newmvar\ma{a\or b\or c}
\newcommand\mb{\ma[2]}
\newcommand\mc{\ma[3]}
\newcommand\mctx[1][1]{\BoundV{\ctx[#1]}}
\newmvar\g{g\or h\or k}
\newcommand\mg[1][1]{\BoundV{\g[#1]}}
\newcommand\Schur{\mathsf{Schur}}
\newcommand\schur{\mathsf{schur}}
\newcommand\transformation{\textbf{transformation}}
\newcommand\plane{\textbf{plane}}
\renewcommand\amalg{+}

\newcommand\keys{\mathbb K}
\newcommand\cmp{\mathsf{compare}}
\newcommand\LT{\mathrm{\DConst{Lt}}}
\newcommand\GT{\mathrm{\DConst{Gt}}}
\newcommand\EQ{\mathrm{\DConst{Eq}}}
\newcommand\Rel{\mathop{\mathsf{Rel}}\nolimits}
\newcommand\Sandwich[1]{#1_{\TConst\bot}^{\TConst\top}}
\newcommand\sandwich[1]{\mathrel{#1_{\bot}^{\top}}}
\newcommand\pminfty{\DConst{\pm\infty}}
\newcommand\mpinfty{\DConst{\mp\infty}}
\newcommand\pinfty{\DConst{\infty}}
\newcommand\minfty{\DConst{-\infty}}
\newcommand\trichotomy{\mathop{\mathsf{trich}}\nolimits}
\newcommand\interval[1][\keys]{\mathsf{interval}_{#1}}
\newcommand\BST[3][\parent]{\mathbb T_{#1{#2}}{#3}}
\newcommand\TLeaf[1][]{\subimplicit[#1]{\mathop\Leaf}}
\newcommand\TNode[1][\triple]{\Node#1}
\newcommand\Btx{\Ctx_{\BST[]{}{}}}
\newcommand\structBST{\mathbf{BST}}
\newcommand\tleaf{\mathsf{leaf}}
\newcommand\tnode{\mathsf{node}}
\newcommand\gast[1][\G]{#1\mathbin{\TConst\boxastSymb}}
\newcommand\lollybox[1][\G]{#1\mathrel{\TConst{\relbar\joinrel\mathrel{\boxastSymb}}}}
\newcommand\alg{\mathbf{alg}}
\newcommand\abs[1]{\left\lvert#1\right\rvert}
\newcommand\muast{\boxast}
\newcommand\rec{\mathrm{rec}}
\newcommand\Pair[1]{#1\mathbin{\TConst\times}}
\newcommand\insrtElim{\mathop{\mathsf{insert}}\nolimits^{\mathsf{Elim}}}
\newcommand\insrtCore{\mathop{\mathsf{insert}}\nolimits^{\mathsf{Core}}}

\newcommand\Fam{\Category{Fam}}
\newop{\costrength}{costrength}
\newcommand\RED{\DConst{r}}
\newcommand\BLU{\DConst{b}}
\newcommand\lift{\mathop{\mathsf{\DConst{lift}}}\nolimits}
\newcommand\coerce{\mathop{\mathsf{coerce}}\nolimits}
\newcommand\goLeft{\mathsf{goLeft}}
\newcommand\flg[1][]{\subimplicit[{#1}]{\mathop{\mathbf{flog}}\nolimits}}
\newcommand\subtree{\mathsf{subtree}}
\newcommand\subtreeCore{\mathsf{subtree}^{\mathrm{Core}}}
\newcommand\subtreeAux{\mathsf{subtree}^{\mathrm{aux}}}
\newcommand\subtreeAuxCore{\mathsf{subtree}^{\mathrm{auxCore}}}
\newcommand\canonise[1][]{\mathop{\mathsf{canonise}}\nolimits_{#1}}
